\section{电子束为什么会洗浅单电子的振荡}
\label{chap:feqo-averaging}

式~\eqref{eq:feqo-pinem-sidebands} 是给定轨迹、给定到达时刻的条件概率。电子束内的每一个电子仍可经历幺正相位调制，却不一定经历同一个 \(\beta\)。设横向坐标 \(\mathbf R\) 与到达时刻 \(t_0\) 的联合概率密度 \(f(\mathbf R,t_0)\geq0\)，且 \(\int f\dd^2R\dd t_0=1\)。未逐个记录这两项时，计数比例为
\begin{equation}
\overline P_\ell=\int\dd^2R\dd t_0\,f(\mathbf R,t_0)
J_\ell^2\!\left(2|\beta(\mathbf R,t_0)|\right).
\label{eq:feqo-ensemble-pinem}
\end{equation}
这里先算概率再平均，因为不同电子是不同制备事件；\(J_\ell^2(2|\mathbb E\beta|)\) 通常不等于式~\eqref{eq:feqo-ensemble-pinem}。

一个可完整算出的例子能显示尺度。设电子束的横向分布为 \(f_R(R)=(2\pi\sigma_R^2)^{-1}e^{-R^2/(2\sigma_R^2)}\)，近场耦合为 \(\beta(R)=\beta_0e^{-R^2/(2L_E^2)}\)。在 \(|\beta_0|\ll1\) 时，\(J_1(2|\beta|)=|\beta|+O(|\beta|^3)\)，所以
\begin{align*}
\overline P_1
&=|\beta_0|^2\int_0^\infty\frac{R\dd R}{\sigma_R^2}
e^{-R^2/(2\sigma_R^2)-R^2/L_E^2}+O(|\beta_0|^4)\\
&=\frac{|\beta_0|^2}{1+2\sigma_R^2/L_E^2}+O(|\beta_0|^4).
\end{align*}
积分用了 \(\int_0^\infty Re^{-aR^2}\dd R=(2a)^{-1}\)。束斑宽度 \(\sigma_R\) 与场尺度 \(L_E\) 的比值直接控制可见边带；若 \(\sigma_R\gtrsim L_E\)，用中心电子的 \(\beta_0\) 拟合整个束流会系统性高估耦合均匀性。强场时仍须计算原始 Bessel 积分，不能延用这个弱场表达。

光脉冲与电子到达时间的平均是另一件事。若近场包络 \(\beta(t_0)=\beta_0e^{-t_0^2/(2\tau_L^2)}\)，到达时刻服从宽度 \(\sigma_t\) 的零均值高斯分布，同样在弱耦合下
\[
\overline P_1=\frac{|\beta_0|^2}{\sqrt{1+2\sigma_t^2/\tau_L^2}}+O(|\beta_0|^4).
\]
空间平均含二维 Jacobian，时间平均只含一维积分，两者的衰减因子不同。若还有探测器能量响应 \(R(y|E)\)，它在这些理想概率形成后再卷积；不能把所有对比度降低都叫作“能量展宽”。

强度平均与相位平均甚至改变不同的量。设每个电子都遇到同一个 \(|\beta|=b\)，但光学相位 \(\phi\) 随发射事件变化。由式~\eqref{eq:feqo-pinem-sidebands}，任一事件的能谱仍是 \(P_\ell=J_\ell^2(2b)\)，与 \(\phi\) 无关；因而均匀随机相位不会把这些峰直接洗浅。可是若把事件相位丢弃，边带密度矩阵的 \(\ell,\ell'\) 元会乘
\[
\langle e^{\ii(\ell-\ell')\phi}\rangle_\phi.
\]
例如相位在圆周上均匀分布时，该因子对 \(\ell\ne\ell'\) 为零；若相位抖动在一个参考相位附近服从方差 \(\sigma_\phi^2\) 的窄高斯分布，则模长是 \(e^{-(\ell-\ell')^2\sigma_\phi^2/2}\)。因此普通能谱可能几乎不变，后续时间聚束或参考场干涉却消失。要判断“对比度降低”究竟来自强度不均、相位不同步还是探测器分辨率，应比较不同类型的可观测量，不能仅凭一条能谱判断。

单一能谱只依赖 \(|\beta|\)。为了读出 \(\arg\beta\)，可让电子经过相位已知的第二个场。若两次相互作用之间的色散相位可忽略，乘积 \(U(\beta_2)U(\beta_1)=U(\beta_1+\beta_2)\)，其中指数在此平移不变模型里可交换。扫描 \(\beta_2=\beta_{\rm ref}e^{\ii\varphi}\) 后，\(|\beta_1+\beta_2|^2\) 含 \(2\operatorname{Re}(\beta_1\beta_2^*)\)，场相位遂转成人口变化。两区距离较大时必须把中间传播相位插入，简单相加不再成立。

\begin{exercise}
用二维高斯束的积分证明上式的 \(1+2\sigma_R^2/L_E^2\) 分母，并解释为什么在 \(\sigma_R\to0\) 时恢复单电子结果。
\end{exercise}
